\chapter{Strength training}
\index{Strength training}

\medskip
\noindent\textit{Educational reference; always seek professional advice for individual care.}

\section*{Definition and scope}
Strength training, also called resistance training, is any form of exercise in which the muscles work against a force, or resistance, that is heavier than their everyday load. The resistance can come from free weights, weight machines, resistance bands, medicine balls, or simply the body's own weight through exercises such as push-ups, squats, and lunges. The goal is to make the muscles progressively stronger, larger, more enduring, and better able to perform daily activities. Strength training is distinct from aerobic or cardio exercise such as walking, running, and cycling, although the two complement each other, and most balanced exercise programs include both. It is not a sport reserved for athletes or bodybuilders: it is a foundation of health for people of every age and ability, and local physical activity guidelines recommend that adults include muscle-strengthening activities at least two days a week. Anyone can participate, from teenagers building healthy habits, to older people protecting their independence, to people training with a physiotherapist to recover from injury. Strength training can also be adapted for wheelchair users, people with disabilities, and those managing chronic conditions, using seated exercises, bands, and machines. This chapter explains what strength training involves, why it matters for health, how it works in the body, what to expect when starting, the risks to be aware of, and practical ways to begin safely.

\section*{Benefits and purpose}
The benefits of strength training reach far beyond muscle size. It directly improves the ability to carry groceries, lift children, climb stairs, and get up from a chair, which translates into independence and quality of life as people age. Resistance exercise builds bone density and is a key strategy in preventing and managing osteoporosis, because bones, like muscles, respond to load by becoming stronger. It improves blood sugar control by making muscles more efficient at taking up glucose, reduces blood pressure, and improves cholesterol levels, all of which lower the risk of heart disease and type 2 diabetes. In addition, strength training supports weight management by raising resting metabolism, protects joints by strengthening the muscles around them, and improves balance and coordination, which reduces falls in older people. For mental health, it reduces symptoms of anxiety and depression, improves sleep, and boosts confidence and mood. For people living with chronic conditions such as arthritis, diabetes, or heart disease, appropriately supervised strength training is an accepted part of management, and for older people it is one of the most effective measures against frailty and disability. In combination with aerobic exercise, it provides the full package of health benefit recommended in national guidelines. The evidence base is strong: systematic reviews of thousands of participants show that regular resistance exercise lowers the risk of falls, improves mobility in people with arthritis, reduces blood pressure by a small but meaningful amount, and can reverse some of the muscle loss that accompanies ageing, a condition known as sarcopenia. It also appears to protect cognitive function in later life and is increasingly used alongside other treatments in the management of depression.

\section*{How it works}
When a muscle is asked to generate force against a resistance close to its maximum, the individual muscle fibres experience microscopic damage and mechanical strain. Over the following hours and days the body repairs and reinforces the fibres, building slightly more protein into the muscle. This is why muscles are sore in the day or two after an unfamiliar session, a normal response called delayed-onset muscle soreness, and why rest and recovery between sessions are essential. To keep improving, the load must be progressively increased over weeks and months, a principle known as progressive overload. The brain also adapts: early gains in strength come largely from the nervous system learning to coordinate muscles more efficiently, which is why beginners can get stronger before any visible change in muscle size. Resistance training can be organised in many ways. Common approaches include sets and repetitions of an exercise with a certain load, exercises targeting major muscle groups such as the chest, back, legs, shoulders, arms, and core, and variations in speed, range of motion, and equipment. A well-designed program works each major muscle group at least twice a week, allows recovery between sessions, and progresses the weight or difficulty steadily. The choice of load matters: heavier weights with fewer repetitions build maximal strength, moderate loads with more repetitions build muscle size, and lighter loads with higher repetitions improve muscular endurance, but for general health a simple approach of choosing a weight that can be lifted with good form for eight to twelve repetitions works well for most people. Two to four sets of each exercise are typical, with a minute or two of rest between sets.

\section*{What to expect}
Starting strength training can feel daunting, but the experience is predictable and manageable. A beginner typically begins with light loads or body-weight exercises, learning the correct technique before increasing weight. Sessions are often two to three times a week, with a rest day between sessions for the same muscle groups, and each session lasts anywhere from 20 to 60 minutes. Early on, strength can improve within a few weeks, mainly from better coordination, while visible changes in muscle size take eight to twelve weeks of consistent training. Muscle soreness is normal after new or hard sessions and usually settles within a few days; it is different from the sharp pain of an injury and should not be severe enough to stop normal activity. Beginners should expect to breathe out on the effort and breathe in on the release, to use smooth, controlled movements, and to stop if they feel sharp pain, dizziness, or chest discomfort. Working with a physiotherapist, exercise physiologist, or accredited personal trainer, at least initially, helps ensure technique is safe and the program suits the individual's goals, health, and experience. It is also normal to feel unsure about gym equipment; many gyms offer an induction session that shows new members how to use the machines and free weights safely, and community exercise classes and programs run by local councils provide a welcoming, low-cost entry point.

\section*{When to seek professional advice}
Most healthy adults can begin a sensible strength training program without a medical visit, but a general practitioner or accredited exercise professional should be consulted first in certain situations. People with heart disease, high blood pressure that is not well controlled, diabetes, osteoporosis, arthritis, a history of joint or back problems, or any condition affecting the brain or nervous system should seek advice about appropriate exercises and any precautions needed. Pregnant women and people who have recently had surgery, or who take blood-thinning medication, should also check before starting. Anyone who experiences chest pain, palpitations, severe breathlessness, fainting, or unusual swelling during or after exercise should stop and seek medical review. Experienced trainers can modify exercises for individual needs, and a referral to an exercise physiologist is available through the public health insurance chronic disease management plan for people with a chronic condition.

\textbf{Contact your local emergency services for:}
\begin{itemize}
    \item Chest pain, pressure, or discomfort during or shortly after exercise
    \item Fainting, collapse, or loss of consciousness
    \item Severe breathlessness that does not settle with rest
    \item Sudden severe pain, deformity, or inability to bear weight after an exercise
\end{itemize}

\section*{Risks and cautions}
Strength training is safe when performed sensibly, but injuries do occur, almost always because of poor technique, excessive load, or too little recovery. The most common injuries are muscle strains, tendon irritation, and lower back pain, often from lifting with a rounded spine or from attempting weights beyond current capacity. Rushing to add weight too quickly, training the same muscles every day, and neglecting warm-up all increase risk. There are specific cautions for particular groups: people with osteoporosis should avoid heavy spinal loading and deep forward flexion; those with high blood pressure should avoid holding the breath and straining, known as the Valsalva manoeuvre, which can spike blood pressure; and people with diabetes may need to check blood glucose and adjust medication around sessions. Correct form, gradual progression, adequate hydration, and listening to the body are the cornerstones of risk reduction. For most people the risks of strength training are far smaller than the risks of physical inactivity, which is itself a major contributor to chronic disease and premature death.

\section*{Strength training across the lifespan}
Strength training has a role in every stage of life. In children and adolescents it is safe when properly supervised, improves sports performance, builds healthy bones, and supports good movement habits; the old idea that weight training stunts growth has been disproved, and the emphasis should be on learning technique with light loads or body weight. In adults it underpins fitness, body composition, and metabolic health. In older people it is arguably most valuable of all: regular resistance work slows age-related muscle and bone loss, improves balance and gait, reduces the risk and severity of falls, and allows people to continue living independently. Frailty, once thought to be an unavoidable part of ageing, responds to carefully graded strengthening programs, and group classes tailored to seniors are widely available through community health services, physiotherapy practices, and local government. For women, resistance training supports bone density around the menopause, when the protective effect of oestrogen declines, and can be continued safely through pregnancy with appropriate modifications and medical clearance. Whatever the age or stage of life, the same principles apply: start gently, prioritise technique, progress gradually, and seek advice when in doubt.

\section*{Tips and practical advice}
\begin{itemize}
    \item Start with body-weight exercises such as squats, push-ups against a wall, and step-ups, adding bands or weights only as technique improves
    \item Aim for at least two strength sessions a week, with a day's rest between sessions working the same muscles
    \item Learn the technique for each exercise first, using a mirror, video, or an accredited trainer; prioritise quality of movement over the amount of weight
    \item Progress gradually, increasing the load or repetitions by a small amount only when the current level feels comfortably manageable
    \item Warm up with five to ten minutes of light activity and gentle movement of the joints before the main session, and cool down with stretching afterwards
    \item Breathe naturally, exhaling on the effort, and never hold your breath under a load
    \item Include all major muscle groups over the week, balancing pushing, pulling, squatting, hinging, and core work
    \item Let muscles recover: sleep well, eat enough protein and food overall, and allow 48 hours before retraining the same muscle group
    \item Keep a simple record of the exercises, loads, and repetitions so that progress can be tracked and increased safely
    \item Vary exercises every few months to keep the program interesting and to challenge muscles in new ways
    \item Combine strength training with aerobic activity, such as brisk walking, to meet the full national physical activity guidelines
    \item Stop any exercise that causes sharp pain and have persistent pain assessed before resuming
    \item Consider joining a supervised class, which provides structure, motivation, and guidance from an instructor
\end{itemize}

\section*{Sources and further reading}
The following organisations provide reliable information about strength training, physical activity, and exercise safety.

\begin{itemize}
    \item Australian Department of Health and Aged Care. \textit{Physical activity and exercise guidelines for all Australians.} https://www.health.gov.au/
    \item healthdirect Australia. \textit{Resistance and strength training for adults.} https://www.healthdirect.gov.au/
    \item NHS. \textit{Strength and flexibility exercise: how to get started.} https://www.nhs.uk/
    \item World Health Organization (WHO). \textit{Physical activity fact sheet and guidelines.} https://www.who.int/
    \item Mayo Clinic. \textit{Strength training: get stronger, leaner, healthier.} https://www.mayoclinic.org/
    \item Exercise and Sports Science Australia (ESSA). \textit{Position statements and consumer resources on resistance training.} https://www.essa.org.au/
\end{itemize}

\clearpage
